\documentclass[10pt]{article}

\usepackage[latin1]{inputenc}
\usepackage{amsmath, amssymb, amsfonts, amsthm}
\usepackage{upgreek}
\usepackage{amsthm}
\usepackage{fullpage}
\usepackage{graphicx}
\usepackage{cancel}
\usepackage{subfigure}
\usepackage{mathrsfs}
\usepackage{outlines}
\usepackage[font={sf,it}, labelfont={sf,bf}, labelsep=space, belowskip=5pt]{caption}
\usepackage{hyperref}

\usepackage{fancyhdr}
\usepackage[title]{appendix}

\DeclareMathOperator{\sgn}{sgn}

\pagestyle{fancy}
\headheight 24pt
\headsep    12pt
\lhead{Low Dimensional Shape Optimization}
\fancyfoot[C]{} % hide the default page number at the bottom
\lfoot{}
\rfoot{\thepage}
\renewcommand{\headrulewidth}{0.4pt}
\renewcommand\footrulewidth{0.4pt}
\providecommand{\abs}[1]{\lvert#1\rvert}
\providecommand{\norm}[1]{\lVert#1\rVert}
\providecommand{\dx}{\, \mathrm{d}x}
% \providecommand{\vint}[2]{\int_{#1} \! #2 \, \mathrm{d}x}
% \providecommand{\sint}[2]{\int_{\partial #1} \! #2 \, \mathrm{d}A}
\renewcommand{\div}{\nabla \cdot}
\providecommand{\shape}{\Omega(p)}
\providecommand{\boundary}{\partial \shape}
\providecommand{\vint}[1]{\int_{\shape} \! #1 \, \mathrm{d}x}
\providecommand{\sint}[1]{\int_{\boundary} \! #1 \, \mathrm{d}A}
\providecommand{\pder}[2]{\frac{\partial #1}{\partial #2}}
\providecommand{\tder}[2]{\frac{\mathrm{d} #1}{\mathrm{d} #2}}
\providecommand{\evalat}[2]{\left.#1\right|_{#2}}
\newcommand{\defeq}{\vcentcolon=}
\newtheorem{lemma}{Lemma}

\makeatletter
\usepackage{mathtools}
\newcases{mycases}{\quad}{%
  \hfil$\m@th\displaystyle{##}$}{$\m@th\displaystyle{##}$\hfil}{\lbrace}{.}
\makeatother

\title{Low Dimensional Shape Optimization}
\author{Julian Panetta}

% BEGIN DOCUMENT
\begin{document}

\maketitle

We seek to optimize a small set of shape parameters, $p \in
\mathbb{R}^{n_p}$, that define an object $\shape \subset \mathbb{R}^n$.
The objective, $J$, is a volume integral of some function, $j$, of the stress
tensor under a particular loading scenario (using linear elasticity):
\begin{align}
    J(p) = \vint{j(x, \sigma(u)(x))} \label{eqn:objective} \\
    \text{s.t.}\quad \begin{mycases}
        -\div \sigma(u) = f(p)   & \text{in}\, \shape \\
        \sigma(u) \hat{n} = g(p) & \text{on}\, \Gamma_N(p) \\
        u = 0 & \text{on}\, \Gamma_D
        \label{eqn:state}
    \end{mycases}
\end{align}
where $\sigma(u) = C:\epsilon(u),\ \epsilon(u) = \frac{1}{2}\left(\nabla u + (\nabla u)^T\right)$
(i.e. $\sigma_{ij} = C_{ijkl} \epsilon_{kl},\ \epsilon_{kl} = \frac{1}{2}\left(u_{k,l} + u_{l,
k}\right)$), $\hat{n}$ is the outward-pointing normal, and $\Gamma_N(p) \defeq
\boundary \setminus \Gamma_D$. Notice the inclusion of a fixed Dirichlet
boundary; we discuss in Section \ref{sec:no_dirichlet} how to optimize without
this constraint. We focus on the task of computing this objective function's
gradient with respect to the shape parameters, neglecting for now the terms
that penalize shapes differing too greatly from the initial design,
$\Omega(p_0)$.

Using a finite difference-based differentiation scheme is difficult due to the
challenge of choosing a reasonable step size. Not only is a good step size, e.g.
for a rectangle's rotation parameter, not immediately clear, but also as
we shrink the step size to get a more accurate approximation, the
noise/errors in the PDE constraint solution will dominate. Instead, we use the
adjoint method \cite{ambrad:website} to compute the gradient. This will involve
solving a single additional ``adjoint'' PDE of the same type as
the state equation (\ref{eqn:state}) and computing a volume and surface
integral per parameter. The derivation here follows roughly that in
\cite{Allaire2008909}, but it handles volume forces, it allows the entire Neumann
boundary to evolve, and it is phrased in terms of shape parameters.

\section{The Lagrangian and the Adjoint Method}
We want to differentiate (\ref{eqn:objective}) with respect to parameter
$p_i$. However, proceeding in a straight-forward way quickly hits a dead end:
$$
\pder{}{p_i} J(p) = \pder{}{p_i} \vint{j(x, \sigma)} = \sint{j(x, \sigma) (v_i(p) \cdot \hat{n})} + \vint{\pder{}{p_i} j(x, \sigma)},
$$
where $v_i(p)$ is the velocity of the boundary while varying $p_i$. This
formula follows immediately from Reynold's transport theorem, but it also can be
understood intuitively: the integral changes due to volume entering/leaving the
integration domain (the first term) and due to the integrand changing (the
second). The first term is easy to compute, but the second requires determining
how the equilibrium displacement changes as the shape changes (hard). The
adjoint method is a tool to sidestep that computation.

Following the traditional adjoint method derivation \cite{ambrad:website}, we
write the Lagrangian for the PDE-constrained optimization. We introduce
Lagrange multipliers $\lambda\colon \overline{\Omega}(p) \to \mathbb{R}^n$ for the Neumann
constraints (\ref{eqn:state}), but set $\lambda = 0$ on $\Gamma_D$  since we do
not need a term in the Lagrangian for the Dirichlet condition (we will always
have $u$ and all its perturbations vanish on the fixed $\Gamma_D$). The
Lagrangian is thus:
$$
\mathcal{L}(p, u, \lambda) = \vint{j(x, \sigma(u))} - \vint{\lambda \cdot \left[\div \sigma(u) + f(p)\right]}
+ \sint{\lambda \cdot \left[\sigma(u) \hat{n} - g(p)\right]},
$$
where $u = 0$ on $\Gamma_D$ is a displacement field independent of $p$ (i.e. not
necessarily the equilibrium displacement). Notice that if we take $u = u_*(p)$,
the solution to the state equation, then $\mathcal{L}(p, u_*(p), \lambda) =
J(p)$, independent of $\lambda$. Therefore:
$$
\pder{J}{p_i}(p) = \tder{\mathcal{L}(p, u_*(p), \lambda)}{p_i}
= \pder{\mathcal{L}}{p_i}(p, u_*(p), \lambda) + \left<\pder{\mathcal{L}}{u}(p, u_*(p), \lambda), \pder{u_*}{p_i}(p)\right>,
$$
where $\pder{\mathcal{L}}{u}$ is a linear functional that, when fed a vector
field $\phi$, computes the derivative of $\mathcal{L}$ in the $\phi$ direction (i.e. the Frech\'et derivative):
$$
\left<\pder{\mathcal{L}}{u}(p, u_*(p), \lambda), \phi\right> =
\evalat{\pder{\mathcal{L}(p, u_*(p) + h \phi, \lambda)}{h}}{h = 0}.
$$
Here is the key insight of the adjoint method: if we can choose $\lambda = \lambda_*(p)$ so
that this linear functional is zero (vanishes on all admissible input vector fields), then
we don't have to compute the problematic term $\pder{u_*}{p_i}(p)$. Then we
just have:
\begin{equation}
\pder{J}{p_i}(p) = \pder{\mathcal{L}}{p_i}(p, u_*(p), \lambda_*(p)).
\label{eqn:gradient_format}
\end{equation}
Finding such a $\lambda_*(p)$ is called the adjoint problem.

\section{The Adjoint Problem}
Soon we will derive the adjoint equations by setting
$$
\left<\pder{\mathcal{L}}{u}(p, u_*(p), \lambda), \phi \right> \overset{!}{=} 0 \quad \forall \phi,
$$
but first we rewrite $\mathcal{L}$ in a more convenient form by restating
in terms of strains and integrating by parts:
\begin{align}
    \mathcal{L}(p, u, \lambda) = \vint{j(x, \sigma(u))} &- \vint{\lambda \cdot \left(\div \left[ C:\epsilon(u)\right] + f(p)\right)}
    + \sint{\lambda \cdot \left(\left[C : \epsilon(u)\right] \hat{n} - g(p)\right)} \notag \\
    = \vint{j(x, \sigma(u))} &+ \overbrace{\vint{\epsilon(\lambda) : C : \epsilon(u)}
- \cancel{\sint{\lambda \cdot \left[C:\epsilon(u)\right]\hat{n}}}}^{\text{from I.B.P. identity (\ref{eqn:IBP_1}) in Appendix \ref{sec:IBP}}} \notag \\
        &+ \cancel{\sint{\lambda \cdot \left[C:\epsilon(u)\right]\hat{n}}}
         - \vint{\lambda \cdot f(p)} - \sint{\lambda \cdot g(p)} \notag \\
    = \vint{j(x, \sigma(u))} &+ \vint{\epsilon(\lambda) : C : \epsilon(u)} - \vint{\lambda \cdot f(p)} - \sint{\lambda \cdot g(p)}
    \label{eqn:lagrangian}
\end{align}

Now we're ready to differentiate with respect to $u$. Let $\phi$ be an
arbitrary perturbation vanishing on $\Gamma_D$. Then:
\begin{align*}
    \left<\pder{\mathcal{L}}{u}(p, u, \lambda), \phi \right> &= \evalat{\tder{}{h}}{h = 0}
\left(\vint{j\big(x, \sigma(u + h\phi)\big) + \epsilon(\lambda) : C : \epsilon(u + h\phi)}\right) \\
&= \vint{\evalat{\pder{}{h}}{h = 0}j\big(x, \sigma(u) + h \sigma(\phi)\big) + \evalat{\pder{}{h}}{h = 0}\epsilon(\lambda) : C : \left[\epsilon(u) + h\epsilon(\phi)\right]},
\end{align*}
where we used the linearity of $\epsilon$ and $\sigma$. Now, defining $j'$ to be
the rank 2 tensor of partial derivatives of $j$ wrt. $\sigma_{ij}$, we get by the chain rule:
\begin{align*}
    \left<\pder{\mathcal{L}}{u}(p, u, \lambda), \phi \right>
    &= \vint{j'\big(x, \sigma(u)\big) : \sigma(\phi) + \epsilon(\lambda):C:\epsilon(\phi)} \\
    &= \vint{\left[j'\big(x, \sigma(u)\big) + \epsilon(\lambda)\right] : C : \epsilon(\phi)}
\end{align*}
Integrating by parts (see (\ref{eqn:IBP_2}) in Appendix \ref{sec:IBP}) to move the strain operator off $\phi$ we get:
\begin{align*}
    \left<\pder{\mathcal{L}}{u}(p, u, \lambda), \phi \right> &=
    \vint{\left(\div C : \left[j'\big(x, \sigma(u)\big)+ \epsilon(\lambda)\right]\right) \cdot \phi}
    + \sint{\left(C : \left[j'\big(x, \sigma(u)\big)+ \epsilon(\lambda)\right]\right) \hat{n} \cdot \phi}
\end{align*}
Finally, we derive the adjoint equation by plugging in $u_*(p)$ and requiring
this derivative to be zero:
$$
\vint{\left(\div C : \left[j'\Big(x, \sigma\big(u_*(p)\big)\Big)+ \epsilon(\lambda)\right]\right) \cdot \phi}
+ \sint{\left(C : \left[j'\Big(x, \sigma\big(u_*(p)\big)\Big)+ \epsilon(\lambda)\right]\right) \hat{n} \cdot \phi}
\overset{!}{=} 0.
$$
For this to be true for all $\phi = 0$ on $\Gamma_D$, we must have:
\begin{align}
    \begin{mycases}
        -\div \sigma(\lambda) = \div \left[C: j'\Big(x, \sigma\big(u_*(p)\big)\Big)\right] & \text{in}\, \shape \\
        \sigma(\lambda) \hat{n} = \left[C : j'\Big(x, \sigma\big(u_*(p)\big)\Big)\right] \hat{n} & \text{on}\, \Gamma_N(p) \\
        \lambda = 0 & \text{on}\, \Gamma_D
    \end{mycases}
\end{align}
where the Dirichlet condition comes from the requirement that $\lambda$ vanish on $\Gamma_D$.
This is a well-posed elastostatic problem whose solution is called the ``adjoint
state,'' $\lambda_*(p)$.

\section{The Gradient}
We compute the right-hand side of (\ref{eqn:gradient_format}) by applying Reynold's
transport theorem and (\ref{eqn:surface_derivative}) from Appendix
\ref{sec:der_sint} to (\ref{eqn:lagrangian}):
\begin{align*}
    \pder{\mathcal{L}}{p_i} = &\sint{\left[j\big(x, \sigma(u)\big) + \epsilon(\lambda) : C : \epsilon(u) -
  \lambda \cdot f(p) - H \lambda \cdot g(p) - \pder{\big(\lambda \cdot g(p)\big)}{\hat{n}}\right] (v_i(p) \cdot \hat{n})} \\
  - &\sint{\lambda \cdot \pder{g(p)}{p_i}}
- \vint{\lambda \cdot \pder{f(p)}{p_i}},
\end{align*}
where $H = \div\hat{n}$ is the curvature (twice the mean curvature in 3D).

Now, using (\ref{eqn:gradient_format}) and plugging in equilibrium and adjoint states $u_*(p)$ and $\lambda_*(p)$, we get:
\begin{align}
    \begin{split}
        \pder{J}{p_i}(p) &= \sint{\left[j\big(x, \sigma(u_*)\big) + \epsilon(\lambda_*) : C : \epsilon(u_*) -
    \lambda_* \cdot f(p) - H \lambda_* \cdot g(p) - \pder{\big(\lambda_* \cdot g(p)\big)}{\hat{n}}\right] (v_i(p) \cdot \hat{n})} \\
   &-\sint{\lambda_* \cdot \pder{g(p)}{p_i}}
    - \vint{\lambda_* \cdot \pder{f(p)}{p_i}}.
    \end{split}
    \label{eqn:partials}
\end{align}
This equation tells us how to compute each component of the gradient provided we have:

\begin{enumerate}
    \item (Normal) boundary velocities induced by each parameter.
    \item How $f$ and $g$ change with $p_i$.
    \item How $\lambda_* \cdot g(p)$ changes in the normal direction.
\end{enumerate}
Item 1 is relatively easy and will be discussed in the following subsections.
Items 2 and 3 are more problematic. In particular, item 3 is not well-defined
as stated since $g$ is only defined on the boundary. As discussed in Appendix
\ref{sec:der_sint}, the {\it existence} of an extension into a neighborhood of
the boundary is needed at the very least. As discussed there, it may be
possible to drop the normal derivative term---the case here is less clear since
$\lambda_*$ already is defined inside the object. Another possible way to
side-step this term is to blur all surface tractions into volume forces and set
$g = 0$ (which is how the current meshfree discretization works anyway).
However, we'd still have to figure out how to differentiate the resulting
volume force, $f$, as mentioned in item 2.

\subsection{Boundary Velocities for CSG Trees}
\label{sub:For CSG Trees}
In the case of objects represented by CSG trees, $p$ is a vector of all the
parameters appearing in any primitive of the tree. As claimed in
\cite{Chen:2007}, assuming primitives in a CSG object intersect
transversely, the boundary can be decomposed into segments controlled by a
single primitive's design parameters. Then, $v_i(p)$ is zero on every
boundary segment not belonging to $p_i$'s primitive(s). On each boundary
segment where it is nonzero, $v_i(p)$ can be found by analyzing the
corresponding primitive in isolation.

The boundary velocities for translation parameters are trivial. Rotation is not
so hard either: for a rotation around axis $\hat{a}$ through point $c$, the
velocity at a boundary point $b$ is $\hat{a} \times (b - c)$. Size parameters
for some primitives are easy (boxes), but for others they will require more
thought (ellipsoids).

Incidentally, it is unclear how \cite{Chen:2007} maintains its fixed Dirichlet
boundary. It must only consider examples where there are no parameters affecting
the Dirichlet boundary. In our case this is too restrictive, so we will need to
relax the requirement of a nonempty Dirichlet boundary (see Section
\ref{sec:no_dirichlet}).

\subsection{Boundary Velocities for Parametrized Level Set Functions} % (fold)
\label{sub:Gradients for Arbitrary Level Set Functions}
Suppose now that we do not have the structure of a CSG tree and instead have
only a parameterized level set function $\Phi(p, x)$ defining the object:
$$
    \begin{cases}
        \Phi(p, x) > 0 &\text{ for } x \in \shape \\
        \Phi(p, x) = 0 &\text{ for } x \in \boundary \\
        \Phi(p, x) < 0 &\text{ for } x \in \mathbb{R}^n \setminus \overline{\Omega}(p) \\
    \end{cases}
$$
We can find the boundary velocity by considering the position of a particular point
on the boundary as a function, $x(p)$. The boundary velocity is, by definition, $v_i(p)(x)
= \pder{x}{p_i}$, and the chain rule gives:
$$
\Phi(p, x(p)) = 0 \quad \Longrightarrow \quad \pder{\Phi(p, x(p))}{p_i} = 0 =
\pder{\Phi}{p_i}(p, x(p)) + \nabla \Phi \cdot v_i(p)\
$$
Recalling the basic level set result, $\hat{n} = \frac{\nabla \Phi}{\norm{\nabla \Phi}}$:
$$
\norm{\nabla \Phi} \hat{n} \cdot v_i(p) = -\pder{\Phi}{p_i}(p, x(p))  \quad \Longrightarrow \quad
\hat{n} \cdot v_i(p) = \frac{-1}{\norm{\nabla \Phi}} \pder{\Phi}{p_i}(p, x(p)).
$$
This equation gives us the normal velocity for points on any boundary of the object.

\section{Getting Rid of $\Gamma_D$}
\label{sec:no_dirichlet}
The fixed Dirichlet boundary $\Gamma_D$ is incompatible with our goals.
Unfortunately, all of the literature seems to rely on this boundary condition.
Ideally, we would instead use the no-rigid-motion constraints from our
Worst-case Structural Analysis paper and allow the entire boundary to evolve:
\begin{align}
    J(p) = \vint{j(x, \sigma(u)(x))} \notag \\
    \text{s.t.}\quad \begin{mycases}
        -\div \sigma(u) = f(p)   & \text{in}\, \shape \\
        \sigma(u) \hat{n} = g(p) & \text{on}\, \boundary \\
        \vint{u} = 0 , & \vint{u \times x} = 0
    \end{mycases}
    \label{eqn:state_nodirichlet}
\end{align}
Introducing the Lagrange multiplier {\it vectors}, $\lambda_t$ and $\lambda_r$,
for the zero net translation and zero net rotation constraints respectively,
the Lagrangian becomes:
\begin{align*}
\mathcal{L}(p, u, \lambda, \lambda_t, \lambda_r) = &\vint{j(x, \sigma(u))} - \vint{\lambda \cdot \left[\div \sigma(u) + f(p)\right]}
+ \sint{\lambda \cdot \left[\sigma(u) \hat{n} - g(p)\right]} + \\ 
&\lambda_t \cdot \vint{u} + \lambda_r \cdot \vint{u \times x}
\end{align*}
Again, $\mathcal{L}$ is independent of $\lambda$, $\lambda_t$, $\lambda_r$ when
$u$ is chosen as a solution to (\ref{eqn:state_nodirichlet}). This allows us to
still choose choose the Lagrange multipliers so that:
\begin{align*}
    \left<\pder{\mathcal{L}}{u}(p, u, \lambda, \lambda_t, \lambda_r), \phi \right> &= 0,
\end{align*}
which can be shown to happen when $\lambda_t = 0$, $\lambda_r = 0$, and $\lambda$ satisfies
\begin{align}
    \begin{mycases}
        -\div \sigma(\lambda) = \div \left[C: j'\Big(x, \sigma\big(u_*(p)\big)\Big)\right] & \text{in}\, \shape \\
        \sigma(\lambda) \hat{n} = \left[C : j'\Big(x, \sigma\big(u_*(p)\big)\Big)\right] \hat{n} & \text{on}\, \boundary \\
    \end{mycases}
\end{align}
This, however, is {\it not} a well-posed elastostatic problem: $\lambda_*$ is
determined only up to a rigid motion offset. While this offset doesn't affect the
terms of the gradient where the adjoint state enters though the strain
operator, it does affect the terms like $\lambda_* \cdot f$. More thought is
needed to determine how to proceed.

\section{Contributions}
\label{sec:Contributions}
So far, if we solve the remaining problems, the contributions would be:
\begin{itemize}
    \item Minimize stress-based objectives {\it with respect to shape
        parameters} (ideally, parameters more general than those controlling
        a single primitive).
    \item Optimize without Dirichlet constraints (with no-rigid-motion constraints).
    \item Handle design-dependent loads. There isn't so much literature on this.
        The closest thing I've seen is handling constant pressure Neumann
        conditions on the evolving boundary, but more of a literature search is
        needed.
    \item Penalize visual dissimilarity instead of the traditional volume
          penalties/constraints.
\end{itemize}

\section{What is Needed/Concerns}
So far, it appears the additional derivation/code needed to compute gradients of
stress-based objectives is:
\begin{itemize}
    \item Computing boundary velocity per primitive.
    \item Computing the divergence of the tensor field
          $C: j'\Big(x, \sigma\big(u_*(p)\big)\Big)$. This is likely to have numerical
          robustness/accuracy issues due to cancellation.
    \item Predicting how volume forces or surface tractions change with boundary movement.
\end{itemize}

The additional theory questions/points that need to be addressed are:
\begin{itemize}
    \item How do we reformulate with no-rigid-motion constrains instead of
          Dirichlet boundary conditions while keeping the adjoint problem well-posed?
    \item Can we drop the normal derivative term appearing in (\ref{eqn:partials})?
    \item How do we model loading scenarios changing as the boundary changes?
    \item Find a more rigorous derivation for Appendix \ref{sec:der_sint} (e.g., see footnote).
\end{itemize}

Finally, it is worrying that the gradient specified by (\ref{eqn:partials}) has
the term $\sint{j\big(x, \sigma(u_*)\big) v_i(p) \cdot \hat{n}}$, which encourages
removal of high-stress material. While this term makes sense in context of the
formal derivation, intuitively it is counterproductive; removing high-stress
regions only redistributes their stresses inward and will in fact weaken the
object further. The term must be compensated for by the rest of the gradient's
terms (involving the adjoint state), but it makes me worry that the stress
field doesn't satisfy some of the differentiability properties we implicitly
assume while carrying out the formal computations. This could possibly
invalidate the derivations.

\begin{appendices}
    \section{Integration By Parts Formulas}
    \label{sec:IBP}
    Here we derive the two integration by parts formulas used in this writeup. The
    derivations are straightforward once broken down to their components using
    index notation. First, we derive the one for:
    $$
    \vint{\lambda \cdot \left(\div \left[ C:\epsilon(u)\right]\right)} =
    \vint{\lambda_i \pder{}{x_j}\big[C_{ijkl} \epsilon_{kl}\big]}
    $$
    From product rule:
    $$
    \vint{\lambda_i \pder{}{x_j}C_{ijkl} \epsilon_{kl}(u)} =
    \vint{\pder{}{x_j} \big[\lambda_i C_{ijkl} \epsilon_{kl}(u)\big] - \pder{\lambda_i}{x_j} C_{ijkl} \epsilon_{kl}(u)}
    $$
    Notice that $\pder{\lambda_i}{x_j} C_{ijkl} = \epsilon_{ij}(\lambda)$ by
    the symmetry $C_{ijkl} = C_{jikl}$. Applying Gauss' theorem to the term on the left, we get:
    $$
    \sint{\lambda_i C_{ijkl} \epsilon_{kl}(u)\hat{n}_j} - \vint{\epsilon_{ij}(\lambda)C_{ijkl} \epsilon_{kl}(u)}
    $$
    Reinterpreting these in coordinate-free notation, we arrive at the desired expression,
    \begin{align}
    \sint{\lambda \cdot \big[C : \epsilon(u) \big]\hat{n}} - \vint{\epsilon(\lambda) : C : \epsilon(u)}.
    \label{eqn:IBP_1}
    \end{align}

    Next, we derive the identity used to move the strain operator off $\phi$ in
    $$
    \vint{\left[j'(x, \sigma(u)) + \epsilon(\lambda)\right] : C : \epsilon(\phi)}
    $$
    To simplify notation, let $\gamma(x, u, \lambda)$ be the rank 2 tensor
    $j'(x, \sigma(u)) + \epsilon(\lambda)$. We want to show:
    \begin{align}
    \vint{\gamma(x, u, \lambda) : C : \epsilon(\phi)} =
    -\vint{\left(\div \big[C : \gamma(x, u, \lambda)\big]\right) \cdot \phi}
    + \sint{\big[C : \gamma(x, u, \lambda) \big] \hat{n} \cdot \phi}
    \label{eqn:IBP_2}
    \end{align}
    In component notation we have, by symmetry of constant $C_{ijkl} = C_{ijlk}$,
    $$
    \vint{\gamma_{ij} C_{ijkl} \epsilon_{kl}(\phi)} = \vint{\gamma_{ij} C_{ijkl} \pder{\phi_k}{x_l}} = \vint{\gamma_{ij} \pder{C_{ijkl} \phi_k}{x_l}}
    $$
    By product rule the expression becomes:
    $$
    \vint{\pder{}{x_l}\big[\gamma_{ij} C_{ijkl} \phi_k \big] - \pder{\gamma_{ij}}{x_l} C_{ijkl} \phi_k}.
    $$
    Applying Gauss' theorem to the left term, we get:
    $$
    \sint{\gamma_{ij} C_{ijkl} \phi_k \hat{n}_l} - \vint{\pder{\gamma_{ij}}{x_l} C_{ijkl} \phi_k}.
    $$
    Using the major symmetry $C_{ijkl} = C_{klij}$ and the fact $C_{ijkl}$ is constant, this is:
    $$
    \sint{\gamma_{kl} C_{ijkl} \phi_i \hat{n}_l} - \vint{\pder{\gamma_{kl}C_{ijkl} }{x_j} C_{ijkl} \phi_i}.
    $$
    Reinterpreting in coordinate-free notation, gives the desired formula, (\ref{eqn:IBP_2}).
    \section{Derivative of a Surface Integral}
    \label{sec:der_sint}
    Here we compute the derivative of a boundary integral. Unlike the volume
    integral case, to which Reynold's transport theorem applies, there doesn't
    seem to be a classical theorem available. We derive a slightly more general
    statement than those appearing without proof in many of the topology
    optimization papers (e.g. \cite{Allaire2008909}). This more general version
    has a chance of handling design-dependent loads, but there are still
    problems to be solved as we will soon see.

    The task is to compute:
    $$
    \pder{}{p_i} \sint{h(p, x)}
    $$
    We use Gauss' theorem to rewrite as a volume integral so that we can apply
    Reynold's transport theorem:
    \begin{align*}
        \pder{}{p_i} \sint{h(p, x)} = \pder{}{p_i} \sint{\tilde{h}(p, x)\tilde{n} \cdot \hat{n}}
    = \pder{}{p_i} \vint{\div \big[\tilde{h}(p, x)\tilde{n} \big]}
    \end{align*}

    Here $\tilde{h}$ and $\tilde{n}$ represent an extension of the surface
    fields $h$ and $\hat{n}$ into the volume ($\tilde{h} = h$ and $\tilde{n} =
    \hat{n}$ on $\boundary$). It will turn out that $\tilde{n}$
    is only ever evaluated on the boundary, so only its existence is needed.
    However, {\it we will end up needing to compute normal derivatives of
    $\tilde{h}$}. In \cite{chen2007shape} a similar extension is needed, and
    they discuss using ``transfinite interpolation'' with approximate distance
    fields.
    
    Applying Reynold's transport theorem the derivative becomes:
    $$
    \sint{\div \big[\tilde{h}(p, x)\tilde{n} \big] \big( v_i(p) \cdot \hat{n} \big)}
+ \vint{\pder{}{p_i} \div \big[\tilde{h}(p, x)\tilde{n} \big]}
    $$
    Applying the product rule $\div \big[\tilde{h}(p, x)\tilde{n} \big] =
    \nabla \tilde{h}(p, x) \cdot \tilde{n} + \tilde{h}(p, x)\div \tilde{n}$ on
    the left and exchanging the order of differentiation on the right:
    $$
    \sint{\left[\pder{\tilde{h}(p, n)}{\hat{n}} + h(p, x) H \right]\big(v_i(p) \cdot \hat{n}\big)}
    + \vint{\div \left[\pder{\tilde{h}}{p_i}(p, x)\tilde{n} + \tilde{h}(p, x) \pder{\tilde{n}}{p_i} \right]},
    $$
    where we replaced tilde quantities evaluated on the boundary with there
    pre-extension counterparts and set $H = \evalat{\div \tilde{n}}{\boundary}$,
    which should be the trace of the shape operator (the sum of the principle
    curvatures) for any reasonable extension $\tilde{n}$. We claim the
    boundary normal and its extension change negligibly with $p_i$ so we can
    drop the last term\footnote{I do not have a proof for this, but it seems
    sensible...}. Applying Gauss' theorem:
    $$
    \sint{\left[\pder{\tilde{h}(p, n)}{\hat{n}} + h(p, x) H \right]\big(v_i(p) \cdot \hat{n}\big)}
    + \sint{\left[\pder{\tilde{h}}{p_i}(p, x)\tilde{n}\right] \cdot \hat{n}},
    $$
    so finally we get the identity:
    \begin{equation}
     \pder{}{p_i} \sint{h(p, x)} =
        \sint{\left[\pder{\tilde{h}(p, n)}{\hat{n}} + h(p, x) H \right]\big(v_i(p) \cdot \hat{n}\big)}
        + \sint{\pder{h}{p_i}(p, x)}.
        \label{eqn:surface_derivative}
    \end{equation}
    Though this requires a careful handling of $\tilde{h}$, each term has an
    intuitive meaning: the first accounts for the boundary moving to encounter
    new scalar field values, the second handles dilation/contraction of the
    boundary (since $H v_i(p) \cdot \hat{n}$ gives the change in surface area due to
    normal motion), and the third handles the parameterized changes of $h$ (in
    our case, the changes in the design-dependent load).

    If we can define $\tilde{h}$ so that $\evalat{\pder{\tilde{h}(p, n)}{\hat{n}}}{\boundary} = 0$
    (the most natural choice, which should be possible to do in a
    neighborhood though I don't have a proof), then the first term
    vanishes and we are left with:
    \begin{equation}
     \pder{}{p_i} \sint{h(p, x)} =
     \sint{h(p, x) H * \big(v_i(p) \cdot \hat{n}\big)}
        + \sint{\pder{h}{p_i}(p, x)}.
        \label{eqn:surface_derivative_simple}
    \end{equation}
\end{appendices}
\bibliographystyle{plain}
\bibliography{References}

\end{document}
